\documentclass{article}
\usepackage[utf8]{inputenc}
\usepackage[T1]{fontenc}
\usepackage[icelandic]{babel}
\usepackage[hidelinks]{hyperref}
\usepackage{cleveref}
\usepackage{cleveref-icelandic}
\usepackage{booktabs}
\usepackage{listings}
\usepackage{xcolor}

\lstset{
  language=[LaTeX]TeX,
  basicstyle=\ttfamily\small,
  keywordstyle=\color{blue},
  commentstyle=\color{gray},
  breaklines=true,
  frame=single,
  literate=
    {á}{{\'a}}1 {é}{{\'e}}1 {í}{{\'i}}1 {ó}{{\'o}}1 {ú}{{\'u}}1 {ý}{{\'y}}1
    {Á}{{\'A}}1 {É}{{\'E}}1 {Í}{{\'I}}1 {Ó}{{\'O}}1 {Ú}{{\'U}}1 {Ý}{{\'Y}}1
    {ð}{{\dh}}1 {Ð}{{\DH}}1 {þ}{{\th}}1 {Þ}{{\TH}}1
    {ö}{{\"o}}1 {Ö}{{\"O}}1 {æ}{{\ae}}1 {Æ}{{\AE}}1
    {→}{{$\rightarrow$}}1
}

\title{\texttt{cleveref-icelandic} --- Íslensk beyging í krossvísunum\\
  \large v0.1}
\author{Sindri Smárason}
\date{1. apríl 2026}

\begin{document}
\maketitle

\begin{abstract}
  Pakkinn \texttt{cleveref-icelandic} bætir við íslenskri beygingu í krossvísunum
  sem notaðar eru með \texttt{cleveref}. Með honum er hægt að vísa í
  myndir, töflur, jöfnur og kafla í öllum föllum: nefnifalli, þolfalli,
  þágufalli og eignarfalli, bæði í eintölu og fleirtölu, með sjálfvirkri
  tengingu á helstu forsetningar.
\end{abstract}

\tableofcontents
\newpage

%%% -----------------------------------------------------------------------
\section{Inngangur}

Pakkinn \texttt{cleveref} er nytsamlegur til krossvísana í LaTeX, en hann
styður ekki íslensku og getur aðeins framleitt nefnifall. Í íslensku
þarf að beygja heiti tilvísunar (s.d. mynd, tafla, jafna, kafli) í samræmi
við setningafræðilegt samhengi. \texttt{cleveref-icelandic} leysir þetta vandamál
með því að geyma beygingartöflur fyrir hvert tilvísunarflokkur og leyfa
notandanum að tilgreina fall með einfaldri valkvæðri frástöðu.

%%% -----------------------------------------------------------------------
\section{Kröfur og hleðsla}

\texttt{cleveref-icelandic} krefst eftirfarandi pakka, sem hlaðast sjálfkrafa:

\begin{itemize}
  \item \texttt{hyperref}
  \item \texttt{cleveref}
  \item \texttt{xparse}
\end{itemize}

Pakkinn verður að hlaðast í réttri röð í formálanum:

\begin{lstlisting}
\usepackage[hidelinks]{hyperref}   % einnig colorlinks=true
\usepackage{cleveref}
\usepackage{cleveref-icelandic}
\end{lstlisting}

%%% -----------------------------------------------------------------------
\section{Aðalskipanir}

\subsection{\texttt{\textbackslash icref}}

\begin{lstlisting}
\icref[<fall>]{<merki>}
\end{lstlisting}

Framleiðir beygt heiti tilvísunar ásamt númerinu, með tengingu á
númerið. Valkvæða frumbreytan \texttt{<fall>} getur verið:

\begin{center}
\begin{tabular}{lll}
  \toprule
  Inntak & Fall & Dæmi \\
  \midrule
  (ekkert) & nefnifall & \verb|\icref{fig:a}| \\
  \texttt{nom} eða \texttt{nf} eða \texttt{0} & nefnifall & \verb|\icref[nom]{fig:a}| \\
  \texttt{acc} eða \texttt{þf} eða \texttt{1} & þolfall & \verb|\icref[acc]{fig:a}| \\
  \texttt{dat} eða \texttt{þgf} eða \texttt{2} & þágufall & \verb|\icref[dat]{fig:a}| \\
  \texttt{gen} eða \texttt{ef} eða \texttt{3} & eignarfall & \verb|\icref[gen]{fig:a}| \\
  forsetning & sjálfvirkt fall & \verb|\icref[í]{fig:a}| \\
  \bottomrule
\end{tabular}
\end{center}

Þegar forsetning er gefin upp framleiðir skipunin forsetninguna sjálfa
ásamt beygtri tilvísun:

\begin{lstlisting}
\icref[í]{fig:a}      % → í mynd 1
\icref[vegna]{fig:a}  % → vegna myndar 1
\icref[um]{eq:a}      % → um jöfnu (1)
\end{lstlisting}

\subsection{\texttt{\textbackslash Icref}}

\begin{lstlisting}
\Icref[<fall>]{<merki>}
\end{lstlisting}

Eins og \verb|\icref| en með stórum staf. Notað í upphafi setningar.

\begin{lstlisting}
\Icref{fig:a}         % → Mynd 1
\Icref[dat]{fig:a}    % → Mynd 1
\end{lstlisting}

\subsection{Fleirtala}

Bæði \verb|\icref| og \verb|\Icref| meðhöndla fleirtölu sjálfkrafa þegar
fleiri en ein vísun er tekin fram:

\begin{lstlisting}
\icref{fig:a,fig:b}        % → myndir 1 og 2
\icref[dat]{fig:a,fig:b}   % → myndum 1 og 2
\end{lstlisting}

Númeraraðir eru þjappaðar niður sjálfkrafa:

\begin{lstlisting}
\icref{eq:a,eq:b,eq:c}     % → jöfnur (1) til (3)
\end{lstlisting}

%%% -----------------------------------------------------------------------
\section{Skilgreining á tilvísunarflokkum}

\subsection{\texttt{\textbackslash DeclareIcrefType}}

\begin{lstlisting}
\DeclareIcrefType{<flokkur>}
  {nef et}{þol et}{þag et}{eig et}
  {nef ft}{þol ft}{þag ft}{eig ft}
\end{lstlisting}

Skilgreinir beygingarformin fyrir tilvísunarflokk. \texttt{<flokkur>}
er innra heiti cleveref — \texttt{figure}, \texttt{table},
\texttt{equation}, \texttt{section} o.s.frv.

\begin{lstlisting}
\DeclareIcrefType{figure}
  {mynd}{mynd}{mynd}{myndar}
  {myndir}{myndir}{myndum}{mynda}
\end{lstlisting}

\subsection{\texttt{\textbackslash DeclareIcrefTypeCapital}}

Eins og \verb|\DeclareIcrefType| en skilgreinir form með stórum staf
fyrir \verb|\Icref|:

\begin{lstlisting}
\DeclareIcrefTypeCapital{figure}
  {Mynd}{Mynd}{Mynd}{Myndar}
  {Myndir}{Myndir}{Myndum}{Mynda}
\end{lstlisting}

%%% -----------------------------------------------------------------------
\section{Forsetningartafla}

\subsection{\texttt{\textbackslash DeclareIscrefPreposition}}

Notendur geta bætt við eða breytt forsetningum:

\begin{lstlisting}
\DeclareIscrefPreposition{<forsetning>}{<fall>}
\end{lstlisting}

Þar sem \texttt{<fall>} er \texttt{acc}, \texttt{dat} eða \texttt{gen}.
Til dæmis, til að breyta fallvísun \texttt{á} yfir í þolfall úr þágufalli:

\begin{lstlisting}
\DeclareIscrefPreposition{á}{acc}
\end{lstlisting}

\subsection{Innbyggðar forsetningar}

Eftirfarandi forsetningar eru innbyggðar í pakkann:

\begin{center}
\begin{tabular}{ll}
  \toprule
  Fall & Forsetningar \\
  \midrule
  Þolfall & um, gegnum, kringum, umfram, umhverfis \\
  \midrule
  Þágufall & frá, að, af, andspænis, ásamt, gagnvart, gegn, gegnt, \\
            & handa, hjá, meðfram, nálægt, undan, úr, \\
            & [á, í, fyrir, undir, yfir, við, með, eftir]$^*$ \\
  \midrule
  Eignarfall & til, auk, austan, án, handan, innan, meðal, megin, \\
             & milli, millum, neðan, norðan, ofan, sakir, sunnan, \\
             & sökum, utan, vegna, vestan \\
  \bottomrule
\end{tabular}
\end{center}

\noindent$^*$ Þessar forsetningar geta stýrt þolfalli eða þágufalli eftir
samhengi (hreyfing vs.\ kyrrstöður). Í tæknilegum skjölum er þágufall
sjálfgefið eftir bestu þekkingu höfundar.. Notendur geta
breytt þessu með \verb|\DeclareIscrefPreposition|.

%%% -----------------------------------------------------------------------
\section{Innbyggðir tilvísunarflokkar}
Eftirfarandi flokkar eru innbyggðir með fullum beygingartöflum í eintölu:

\begin{center}
\begin{tabular}{lllll}
  \toprule
  Flokkur & Nf. & Þf. & Þgf. & Ef. \\
  \midrule
  \texttt{figure}   & mynd  & mynd  & mynd  & myndar \\
  \texttt{table}    & tafla & töflu & töflu & töflu  \\
  \texttt{equation} & jafna & jöfnu & jöfnu & jöfnu  \\
  \texttt{section}  & kafli & kafla & kafla & kafla  \\
  \bottomrule
\end{tabular}
\end{center}

\newpage
og í fleirtölu:
\begin{center}
\begin{tabular}{lllll}
  \toprule
  Flokkur & Nf. & Þf. & Þgf. & Ef. \\
  \midrule
  \texttt{figure}   & myndir & myndir & myndum & mynda  \\
  \texttt{table}    & töflur & töflur & töflum & taflna \\
  \texttt{equation} & jöfnur & jöfnur & jöfnum & jafna  \\
  \texttt{section}  & kaflar & kafla  & köflum & kafla  \\
  \bottomrule
\end{tabular}
\end{center}
%%% -----------------------------------------------------------------------
\section{Heildardæmi}

\begin{lstlisting}
\documentclass{article}
\usepackage[icelandic]{babel}
\usepackage[hidelinks]{hyperref}
\usepackage{cleveref}
\usepackage{cleveref-icelandic}

\begin{document}

\begin{figure}
  \caption{Yfirlit}
  \label{fig:yfirlit}
\end{figure}

Niðurstöður eru sýndar í \icref{fig:yfirlit}.
Fjallað er um \icref[acc]{fig:yfirlit} í næsta kafla.
Sjá má frekari upplýsingar \icref[í]{fig:yfirlit}.
\Icref{fig:yfirlit} sýnir samantekt á gögnum.

\end{document}
\end{lstlisting}


%%% -----------------------------------------------------------------------
\section{Þakkir}

Pakkinn byggir á \texttt{cleveref} eftir Toby Cubitt. Hönnun
forsetningakerfisins var innblásið af \texttt{crefthe} og \texttt{xcref}.

\end{document}